\documentclass[a4paper,11pt]{article}
\usepackage{amsmath,amsthm,amsfonts,amssymb,amscd,amstext,vmargin,graphics,graphicx,tabularx,multicol} \usepackage[french]{babel}
\usepackage[utf8]{inputenc}  
\usepackage[T1]{fontenc} 
\usepackage[T1]{fontenc}
\usepackage{amsmath,amssymb}
\usepackage{pstricks-add,tikz,tkz-tab,variations}
\usepackage[autolanguage,np]{numprint} 
\usepackage{enumitem}

\setmarginsrb{1.5cm}{0.5cm}{1cm}{0.5cm}{0cm}{0cm}{0cm}{0cm} %Gauche, haut, droite, haut
\newcounter{numexo}
\newcommand{\exo}[1]{\stepcounter{numexo}\noindent{\bf Exercice~\thenumexo} : \marginpar{\hfill /#1}}
\reversemarginpar


\newcounter{enumtabi}
\newcounter{enumtaba}
\newcommand{\q}{\stepcounter{enumtabi} \theenumtabi.  }
\newcommand{\qa}{\stepcounter{enumtaba} (\alph{enumtaba}) }
\newcommand{\initq}{\setcounter{enumtabi}{0}}
\newcommand{\initqa}{\setcounter{enumtaba}{0}}

\newcommand{\be}{\begin{enumerate}}
\newcommand{\ee}{\end{enumerate}}
\newcommand{\bi}{\begin{itemize}}
\newcommand{\ei}{\end{itemize}}
\newcommand{\bp}{\begin{pspicture*}}
\newcommand{\ep}{\end{pspicture*}}
\newcommand{\bt}{\begin{tabular}}
\newcommand{\et}{\end{tabular}}
\renewcommand{\tabularxcolumn}[1]{>{\centering}m{#1}} %(colonne m{} centrée, au lieu de p par défault) 
\newcommand{\tnl}{\tabularnewline}

\newcommand{\trait}{\noindent \rule{\linewidth}{0.2mm}}
\newcommand{\hs}[1]{\hspace{#1}}
\newcommand{\vs}[1]{\vspace{#1}}

\newcommand{\N}{\mathbb{N}}
\newcommand{\Z}{\mathbb{Z}}
\newcommand{\R}{\mathbb{R}}
\newcommand{\C}{\mathbb{C}}
\newcommand{\Dcal}{\mathcal{D}}
\newcommand{\Ccal}{\mathcal{C}}
\newcommand{\mc}{\mathcal}

\newcommand{\vect}[1]{\overrightarrow{#1}}
\newcommand{\ds}{\displaystyle}
\newcommand{\eq}{\quad \Leftrightarrow \quad}
\newcommand{\vecti}{\vec{\imath}}
\newcommand{\vectj}{\vec{\jmath}}
\newcommand{\Oij}{(O;\vec{\imath}, \vec{\jmath})}
\newcommand{\OIJ}{(O;I,J)}

\newcommand{\bmul}[1]{\begin{multicols}{#1}}
\newcommand{\emul}{\end{multicols}}


\newcommand{\reponse}[1][1]{%
\multido{}{#1}{\makebox[\linewidth]{\rule[0pt]{0pt}{20pt}\dotfill}
}}

\newcommand{\titre}[5] 
% #1: titre #2: haut gauche #3: bas gauche #4: haut droite #5: bas droite
{
\noindent #2 \hfill #4 \\
#3 \hfill #5

\vspace{-1.6cm}

\begin{center}\rule{6cm}{0.5mm}\end{center}
\vspace{0.2cm}
\begin{center}{\large{\textbf{#1}}}\end{center}
\begin{center}\rule{6cm}{0.5mm}\end{center}
}



\begin{document}
\pagestyle{empty}
\titre{Programmer un casse-brique}{}{}{}{5ème}


Ouvrir le logiciel Scratch 2.\\

\textbf{{\large \underline {Étape 1}}} : {\large La balle}




\bmul{2}
\q Supprimer le lutin Scratch "lutin 1" et choisir le lutin "tennis ball" dans la bibliothèque. \\
Renommer ce lutin "Balle" ( clic sur \includegraphics[scale=1]{cassebrique1.eps} )\\

\q La balle de tennis doit rebondir sur les côtés du terrain.\\
Construire le script ci-contre.\\

\columnbreak

\includegraphics[scale=1]{cassebrique2.eps} \\
\emul

\textbf{{\large \underline {Étape 2}}} : {\large La raquette}



\bmul{2}

\initq \q Dessiner un nouveau lutin "Raquette" en forme de rectangle. Penser à centrer votre lutin.\\

\q Voici ci-contre le début du script de ce lutin, quel est son rôle ?\\

\q Prévoir également le déplacement vers la gauche de la raquette et compléter le script.\\

\q Lorsque la balle touche la raquette, elle doit rebondir de 45\degre. Ajouter dans le bon ordre les briques suivantes au script du lutin "Balle" : \\
\columnbreak

\includegraphics[scale=0.8]{cassebrique3.eps} 

\vspace*{1cm}

\includegraphics[scale=0.7]{cassebrique4.eps} \\

\emul

\textbf{{\large \underline {Étape 3}}} : {\large Le mur}\\

\initq \q Dessiner un lutin "Mur" en forme de rectangle allongé en bas de l'écran de jeu.\\



\q Si la balle touche le mur alors la partie est perdue. Programmer cette éventualité à l'aide des blocs suivants :



\begin{center}
 \includegraphics[scale=1]{cassebrique5.eps}\\ 
\end{center}

\q A l'aide de l'instruction \includegraphics[scale=0.8]{cassebrique6.eps}   , placer les lutins "Balle" et "Raquette" dans la même position au début de chaque nouvelle partie.\\

\q Exécutez le jeu et n'oubliez d'enregistrer votre programme. ((Fichier - Enregistrer sous - Bureau - cassebrique.prénom)\\

\textbf{{\large \underline {Étape 3}}} : {\large Les briques}\\

\initq \q C'est le moment de rajouter des briques, pour cela créer un lutin "Brique" de forme rectangulaire. En créer plusieurs soit en dupliquant la première soit en la clonant.\\

\q Lorsque la balle touche les briques, elles doivent disparaître. A vous d'inventer le script ! \\

\textbf{Faîtes évoluer votre jeu en proposant une fin, en ajoutant un score, en mettant des sons aux moments importants ...
}
\end{document}
